\usepackage{amsmath,amssymb,amsthm, amsfonts,mathrsfs}
\usepackage{geometry}
\usepackage{enumerate}
\usepackage{pdfpages}
\usepackage{subcaption}
\usepackage{graphicx}
\usepackage{cite}
\usepackage{tikz-cd} %交换图
\usepackage{tikz}


\usepackage{stix} %symbols and text fonts
\usepackage{charter}
\usepackage[T1]{fontenc}

\usepackage{xcolor}                             % Link color
 \definecolor{custom-blue}{RGB}{0,99,166} 
\usepackage{hyperref}
\hypersetup{
    colorlinks=true, % 启用彩色链接
    linkcolor=purple,  % 普通内部链接的颜色
    citecolor=blue,  % 引用链接的颜色
    filecolor=green,  % 打开本地文件的链接颜色
    urlcolor=red    % 链接到URL的颜色
}

\usepackage{titlesec}  % 用于修改章节标题的格式

% 使用 \titleformat 定义新的 \section 和 \subsection 样式
\titleformat{\section}[block]
  {\normalfont\Large\bfseries}    % 标题字体样式
  {\S\ \thesection}               % 在标题前加上 \S 和节编号
  {1em}                           % 标题和编号之间的间距
  {}                               % 标题后面的内容

\titleformat{\subsection}[block]
  {\normalfont\large\bfseries}    % 子节标题字体样式
  {\S\ \thesubsection}            % 在子节标题前加上 \S 和子节编号
  {1em}                           % 标题和编号之间的间距
  {}                               % 标题后面的内容

  \titleformat{\chapter}[block]
  {\centering\Huge\bfseries}    
  {Chapter \thechapter }           % 居中
  {1em}                           
  {} 

%%%%%%%%页面与页面对齐
\textwidth=6.5in
\textheight=8.7in
\topmargin=0in
\oddsidemargin=0in
\evensidemargin=0in



%%Bold the titles
\makeatletter
\def\th@plain{%
  \thm@notefont{}% same as heading font
  \itshape % body font
}
\def\th@definition{%
  \thm@notefont{}% same as heading font
  \normalfont % body font
}
\makeatother

\newtheorem{theorem}{Theorem}[section]
\newtheorem*{theorem*}{Theorem}
\newtheorem{lemma}[theorem]{Lemma}
\newtheorem{corollary}[theorem]{Corollary}
\newtheorem{proposition}[theorem]{Proposition}
\newtheorem{conjecture}[theorem]{Conjecture}

\theoremstyle{definition}
\newtheorem{definition}[theorem]{Definition}
\newtheorem{example}[theorem]{Example}
\newtheorem{exercise}[theorem]{Exercise}
\newtheorem{problem}[theorem]{Problem}
\newtheorem{remark}{Remark}[section]
\newtheorem{question}{Question}[section]
\renewcommand*{\thequestion}{\Alph{question}}

\DeclareMathOperator{\grad}{grad}
\DeclareMathOperator{\dimension}{dim}
\DeclareMathOperator{\trace}{tr} 
\DeclareMathOperator{\di}{div}
\DeclareMathOperator{\ricci}{Ricci}
\DeclareMathOperator{\dist}{dist}
\DeclareMathOperator{\hess}{Hess}


\DeclareMathOperator{\im}{Im}

\def\zz{{\mathbb Z}}
\def\rr{{\mathbb R}}
\def\cc{{\mathbb C}}
\def\qq{{\mathbb Q}}
\def\nn{{\mathbb N}}
\def\ss{{\mathbb S}}
\def\oo{{\infty}}

\newcommand{\la}{\langle}
\newcommand{\ra}{\rangle}

\newcommand{\dif}{\mathop{}\!\mathrm{d}} % unicode-math 

\newtheoremstyle{coloredtheorem4}
  {\topsep}   
  {\topsep}   
{\normalfont} 
  {}          
  {\color{gray}\bfseries}
  {.}        
  {.5em}      
  {\thmname{#1}\thmnumber{ #2}\thmnote{ (#3)}} 

